\documentclass[11pt,a4paper]{article}
\usepackage[utf8]{inputenc}
\usepackage[T1]{fontenc}
\usepackage{amsmath, amssymb, amsthm}
\usepackage{geometry}
\geometry{margin=2.5cm}
\usepackage{lmodern}
\usepackage[final,expansion=false]{microtype}
\usepackage{booktabs}
\usepackage{array}
\usepackage{enumitem}
\usepackage{xcolor}
\usepackage{titlesec}
\PassOptionsToPackage{hyphens}{url}
\usepackage{hyperref}
\usepackage{xurl}
\hypersetup{
  colorlinks=true,
  linkcolor=blue!50!black,
  urlcolor=blue!50!black,
  citecolor=blue!50!black,
  pdftitle={The Universal Right-Triangle Reflection Sequence -- Companion (conjectures and partial progress)},
  pdfauthor={Vico Bonfioli},
}
\usepackage{parskip}
\setlength{\parindent}{0pt}
\setlength{\parskip}{0.55em}

\titleformat{\section}{\normalfont\Large\bfseries\color{blue!40!black}}{\thesection.}{0.6em}{}

\newtheorem{conjecture}{Conjecture}
\newtheorem{proposition}{Proposition}
\theoremstyle{remark}
\newtheorem{remark}{Remark}
\newtheorem{observation}{Numerical observation}
\theoremstyle{plain}

\title{\textbf{Companion document: partial results, conjectures, and computational evidence}\\[0.3em]
\large Supplement to \emph{The Universal Right-Triangle Reflection Sequence}}
\author{Vico Bonfioli}
\date{\today}

\begin{document}
\maketitle

\paragraph{Preface.}
This companion document collects partial results, conjectures, and
computational evidence supporting the program of the main paper
\texttt{paper.tex}.  \textbf{Nothing here is claimed as a theorem.}
All results are explicitly conditional or empirical: each is labelled
either as a \emph{conjecture} (a target statement we believe but have
not proved), as a \emph{conditional proposition} (a deduction valid
under a hypothesis that is itself only empirically verified at a finite
horizon), or as a \emph{numerical observation} (a finite computational
record).  Cross-references to theorems of the main paper use the same
labels.

\section{The transcendence program for $\beta_2$: conjecture and three attack routes}
\label{sec:transcendence-program}

Define the asymptotic ratio
\[
\beta_2 \;:=\; \lim_{d \to \infty} \frac{u_{d+1}}{u_d}
\]
of the universal right-triangle reflection sequence (existence is
empirical; the latest BFS data through depth $42$ places the ratio in
the narrow window $[1.4718, 1.5019]$ depending on the extrapolation
method).

\begin{conjecture}[Transcendence of $\beta_2$]
\label{conj:beta2-transcendence}
$\beta_2$ is transcendental over $\mathbb{Q}$, and the generating
function $U(t) = \sum_{d \ge 0} u_d t^d$ is transcendental over
$\mathbb{Q}(t)$.
\end{conjecture}

We collect three independent attack routes.  Each rests on a different
finite-horizon hypothesis that is empirically verified and structurally
attackable; the routes are logically independent and would each
individually settle Conjecture~\ref{conj:beta2-transcendence}.

\subsection{Route 1: Christol mod-$p$ via the explicit cyclic-module structure}
\label{ssec:christol-route}

Let $\rho_{\mathrm{abs}}: W \to Q \ltimes M$ be the universal
realisation of the main paper (\S4) with $Q = D_\infty \times \mathbb{Z}/2$
and $M = \mathbb{Z}[Q]/(Y+1, c+1)$.  For each prime $p$, write
\[
S_p \;:=\; \bigl\{\,(\pi(g),\, m \bmod p) \;:\; (g,m) = \rho_{\mathrm{abs}}(w),\ w \in W \,\bigr\}
\;\subset\; Q \times (M/pM).
\]

\begin{proposition}[Mod-$p$ state space is infinite]
\label{prop:modp-state-infinite}
For every prime $p$, $|S_p| = \infty$.  In particular the infinite
family $\{((XY)^n, t_n) : n \ge 1\}$ (with $t_n$ the translation part
of $\rho_{\mathrm{abs}}((R_0 R_2)^n) \bmod p$) consists of pairwise
distinct elements of $Q \times (M/pM)$.
\end{proposition}

\emph{Sketch.}  $XY$ has infinite order in $D_\infty$, so the
$Q$-components $(XY)^n$ are pairwise distinct.  This alone suffices for
infinity of $S_p$ via projection.  The stronger separation on the
$(M/pM)$-coordinate, used in Route~1 below, follows from the witness
monomial $(XY)^{n-1} X \cdot h$ surviving reduction modulo $p$ by the
cyclic-module relations $(Y+1)h = (c+1)h = 0$.

\paragraph{Numerical record.}
An exact-arithmetic Rust BFS in $\mathbb{F}_p[Q]$ enumerates $S_p(d) :=
\{\rho_{\mathrm{abs}}(w) \bmod p : |w|_W \le d\}$ at three primes:
\[
\begin{array}{c|c|c|c}
p & d_{\max} & |S_p(d_{\max})| & \beta_p \text{ (ratio fit, last 10 depths)} \\ \hline
3 & 42 & 65\,428\,964 & 1.4681 \\
5 & 40 & 62\,736\,544 & 1.5048 \\
7 & 40 & 65\,724\,148 & 1.5077
\end{array}
\]
The implementation uses the fixed-degree truncation justified by
$(Y+1)h = (c+1)h = 0$, which bounds the support of any reachable
$m \in M/pM$ by $O(d)$ monomials.

\begin{proposition}[$\beta_2$ transcendence --- Christol route, conditional]
\label{prop:beta2-christol}
Conditional on the Myhill--Nerode separation hypothesis (the states
associated to $\{((XY)^n, t_n)\}_{n \ge 1}$ in the $p$-DFAO for
$u_d \bmod p$ are pairwise distinct future-equivalence classes), $U(t)$
is transcendental over $\mathbb{Q}(t)$ and $\beta_2$ is transcendental
over $\mathbb{Q}$.
\end{proposition}

\emph{Sketch.}  Under the separation hypothesis the $p$-DFAO has
infinitely many states, hence (Allouche--Shallit Thm.~6.6.2) the
$p$-kernel of $U(t) \bmod p$ is infinite, so by Christol 1979 $U(t)
\bmod p$ is not algebraic over $\mathbb{F}_p(t)$.  Adamczewski--Bell
2017 lifts this to non-algebraicity of $U(t)$ over $\mathbb{Q}(t)$, and
Flajolet--Sedgewick (\emph{Analytic Combinatorics} Thm.~VII.7) implies
$\beta_2$ transcendental.

\subsection{Route 2: D-finite exclusion}
\label{ssec:dfinite-route}

\begin{proposition}[No D-finite recurrence at low complexity]
\label{prop:no-dfinite}
On $u_0, \ldots, u_{42}$, no D-finite (P-recursive) recurrence of order
$k \le 9$ with polynomial-in-$d$ coefficients of degree $m \le 7$
exists ($47$ tested $(k, m)$-pairs, exact integer linear algebra
modulo the Mersenne prime $2^{61} - 1$).
\end{proposition}

\begin{proposition}[$\beta_2$ transcendence --- D-finite route, conditional]
\label{prop:beta2-dfinite}
Conditional on extending Proposition~\ref{prop:no-dfinite} from
$(k \le 9, m \le 7)$ to all $(k, m) \in \mathbb{N}^2$, $\beta_2$ is
transcendental over $\mathbb{Q}$.
\end{proposition}

\emph{Sketch.}  By Stanley \emph{Enumerative Combinatorics} Vol.~2
Thm.~6.4.6, every power series algebraic over $\mathbb{Q}(t)$ is
D-finite.  Excluding D-finite at all $(k, m)$ excludes algebraic, hence
(Flajolet--Sedgewick VII.7) transcendence of $\beta_2$.  The natural
structural target for the extension is a Gr\"obner-basis argument on
the skew-polynomial ring $\mathbb{Q}\langle n, S\rangle$ acting against
the cyclic-module relations $(Y+1)h = (c+1)h = 0$.

\subsection{Route 3: Mahler structure}
\label{ssec:mahler-route}

\begin{conjecture}[$U(t)$ is not a class-(1) Mahler function]
\label{conj:not-class1-mahler}
$U(t)$ is not a Mahler function of class (1) in the sense of
Adamczewski--Faverjon 2020.
\end{conjecture}

Were Conjecture~\ref{conj:not-class1-mahler} established, the Adamczewski--Faverjon
framework would yield a third independent route to transcendence of
$\beta_2$, parallel to Propositions~\ref{prop:beta2-christol} and~\ref{prop:beta2-dfinite}.

\subsection{Joint picture}

The three conditionals are logically independent.  For
Conjecture~\ref{conj:beta2-transcendence} to fail, all three hypotheses
would have to fail simultaneously.

\section{Structural exclusions for $U(t)$ at low complexity}
\label{sec:structural-negatives}

The following finite linear-algebra exclusions on $u_0, \ldots, u_{42}$
constrain the form of $U(t)$ at the verifiable horizon.

\begin{proposition}[No low-order linear recurrence]
\label{prop:no-recurrence}
The sequence $u_0, \ldots, u_{38}$ satisfies no linear recurrence over
$\mathbb{Q}$ of order $\le 19$.  (Equivalently: for every $k \le 19$,
the system $u_n = \sum_{j=1}^k c_j u_{n-j}$ has empty rational solution
set on the available terms.)
\end{proposition}

\begin{proposition}[No nonlinear or modular structure at low complexity]
\label{prop:no-recurrence-strong}
On $u_0, \ldots, u_{42}$:
\begin{enumerate}
\item[\textnormal{(i)}] No D-finite recurrence with $k \le 9$ and
$m \le 7$ (Proposition~\ref{prop:no-dfinite}).
\item[\textnormal{(ii)}] No nonlinear polynomial relation $\sum_\alpha q_\alpha(n)\,u_n^{\alpha_0} \cdots u_{n-w}^{\alpha_w} = 0$
with lookback $w \le 5$, total $u$-degree $|\alpha| \le 3$, and
coefficient degree $\le 2$.
\item[\textnormal{(iii)}] No non-trivial modular linear recurrence at
any modulus $m \in \{3, 4, \ldots, 101\}$ except $m = 2$, where the
order-$2$ Fibonacci recurrence $u_d \equiv u_{d-1} + u_{d-2} \pmod 2$
holds throughout the tested range.
\end{enumerate}
\end{proposition}

\begin{proposition}[Coefficient-height gap]
\label{prop:height-gap}
$\log u_d = c \cdot d + o(d)$ with $c \approx 0.404$ (linear regression
on $d \in [25, 42]$, $R^2 \ge 0.999$).  Consequently $u_d$ is not the
value sequence of any $k$-automatic integer sequence (Bell--Coons /
Adamczewski--Bell height-gap theorem) and is not the coefficient
sequence of any Mahler function of bounded complexity.
\end{proposition}

\begin{proposition}[Mod-$2$ algebraicity]
\label{prop:mod2-automaton}
$u_d \bmod 2$ is purely $3$-periodic: $u_d \equiv 1$ for $d \not\equiv 0 \pmod 3$
or $d \le 2$, and $u_d \equiv 0$ otherwise.  Equivalently
$u_d \equiv u_{d-1} + u_{d-2} \pmod 2$ for $d \ge 2$.  Hence
$U(t) \in \mathbb{F}_2[[t]]$ is algebraic over $\mathbb{F}_2(t)$ via a
$4$-state finite automaton.

\emph{Caveat.}  This is a characteristic-$2$ statement and does not
lift to algebraicity over $\mathbb{Q}(t)$.
\end{proposition}

\section{Universality class: scope and known examples}
\label{sec:universality-class}

The mechanism that proves universality in the main paper
(Theorem~\ref{thm:universality}) --- a cyclic-ideal identification
$I_A \cap I_B = \mathbb{Z}[Q]\cdot h$ with explicit central generator
$h$ inside the linear quotient $Q$ of a free product $W = A * B$ ---
admits a natural class of free-product reflection systems.

\begin{conjecture}[Open question: scope of the cyclic-ideal mechanism]
\label{conj:universality-class}
Characterise the pairs $(W, \rho)$ for which the cyclic-ideal mechanism
of Theorem~\ref{thm:universality} applies and yields universality of
the orbit-growth sequence across the family of geometric realisations.
\end{conjecture}

Two examples are currently known.

\paragraph{Example A: right triangles.}  $A = V_4$, $B = \mathbb{Z}/2$,
$Q = D_\infty \times \mathbb{Z}/2$, $h = (c-1)(Y-1)$.  This is the main
paper's Theorem~\ref{thm:universality}.

\paragraph{Example B: hyperbolic two-right-angle quadrilaterals.}
$A = B = V_4$, $W = V_4 * V_4$, $Q = D_\infty \times \mathbb{Z}/2$,
$h = c - 1$ (no $(Y-1)$ factor needed: centrality of $c$ alone suffices
when both factors of the free product contain $c$).

\paragraph{Cases where the mechanism is structurally blocked.}  Acute
and obtuse triangle reflection groups have no central involution
playing the role of $c$.  Equilateral and other triangle groups with
all finite Coxeter labels lack the free-product structure.  The
rectangle group $D_\infty \times D_\infty$ is a direct product, not a
free product.  Right-corner orthoschemes in dimension $n \ge 3$ carry
cross-relations $(R_i R_j)^2 = 1$ between generators of distinct
free-product factors, again destroying the $W = A * B$ structure;
universality in this cyclic-ideal sense is therefore specific to
dimension~2.

\section{Structural reading of universality (algebraic-channel BFS)}
\label{sec:univ-algebraic}

Fix $\rho_{a,b}: W \to \mathrm{Aff}(\mathbb{R}^2)$.  Let $Y$ denote the
linear part of $\rho_{a,b}(R_2)$ and $v_2$ its translation part.  The
identity $Y v_2 = -v_2$ holds for every $(a, b)$ and is the image of a
single canonical word $w_\star \in W$.  Let $N_{\rm alg}$ be the normal
closure of $w_\star$ in $W$.

\begin{observation}[Empirical algebraic form of universality]
\label{obs:alg-structural}
Through depth $22$, on the four test right triangles $(3, 4), (5, 12),
(1, 2), (7, 24)$, the equivalence relation on canonical Coxeter words
generated by $Y v_2 + v_2$ in $\mathbb{Z}[D_\infty]$ exactly recovers
the affine equivalence relation.  Equivalently: $N(T) = N_{\rm alg}$ on
the depth-$\le 22$ filtration for every tested $T$.
\end{observation}

A five-channel parallel BFS verifies this with cumulative state count
$118\,938$ at depth $22$ in every channel.  Reading: the eight
length-$10$ affine relations of Theorem~\ref{thm:symbolic} are not
sporadic length-$10$ coincidences but the first eight consequences of
the single quadratic relation $Y v_2 + v_2 = 0$, made visible by the
canonical-word filtration.

\section{Empirical evidence against finite confluent rewriting}
\label{sec:rewriting}

\paragraph{Knuth--Bendix experiment.}
GAP \texttt{KBMAG} ran Knuth--Bendix completion on $W$ together with
the eight length-$10$ affine relations of Theorem~\ref{thm:symbolic};
the procedure produced $32\,730$ derived rules in approximately $8$
seconds of \texttt{kbprog} computation before hitting the equation-count
cap, never reaching confluence.  This is empirical evidence (not a
proof) that $W/N$ has no finite confluent rewriting system in shortlex
on the standard generators.

\paragraph{Antol\'in--Ciobanu route (structurally blocked).}
The Antol\'in--Ciobanu rationality theorem for conjugacy-growth series
of non-elementary acylindrically hyperbolic groups does not apply to
$U(t)$: the quotient $W / \ker(\rho_{\mathrm{abs}})$ has image inside
$Q \ltimes M \subseteq (\mathbb{Z} \wr \mathbb{Z}) \rtimes V_4$, a
finitely-presented virtually metabelian group, which is not
acylindrically hyperbolic.

\section{Class A and Class B in 3D}
\label{sec:classes-AB}

The 3D right-corner orthoscheme reflection groups fall into three
faithfulness regimes by leg-symmetry, the unconditional Class~C case
being treated in the main paper (Theorem~\ref{thm:master-C-uncond}).

\paragraph{Class A ($a = b = c$, $S_3$-symmetric, polynomial growth).}
The $(1, 1, 1)$ cube-corner orthoscheme is the fundamental chamber of
the affine Weyl group $\widetilde{C}_3$.  For $n \ge 1$,
$u_n = (8 n^2 + c_{n \bmod 5})/5$ with $c = (10, 12, 13, 13, 12)$, with
asymptotic $u_n \sim (8/5) n^2$.  Order-$7$ linear recurrence with
characteristic polynomial $(x-1)^3 \Phi_5(x)$, obtained from the
parabolic Poincar\'e series $W_A(t) = N(t) / ((1-t)^3 \Phi_5(t))$ with
$N(t) = 2 + 3t^2 + 3t^3 + 3t^4 + t^5 + 4t^6$.  This is a clean
restatement of standard affine-Weyl-group material (Bourbaki Ch.~VI~§4;
Humphreys \S4.7) and is recorded here for completeness.

\paragraph{Class B (exactly two equal legs, intermediate exponential rate).}
Through depth $18$ the sequence on $(1, 1, 2)$ is
$1, 4, 9, 18, 35, 67, 129, 249, 480, 925, 1783, 3437, 6625, 12770,
24591, 47298, 90948, 174939, 336544$, with empirical ratio settling
around $\approx 1.92$.  No order-$\le 9$ linear recurrence over
$\mathbb{Q}$ fits the depth-$18$ data.  No closed form is known.

\section{Other deferred material}
\label{sec:deferred}

The following items were deferred from the main paper to this
supplement and remain available in earlier drafts of the project:
\begin{itemize}[leftmargin=*,itemsep=0.15em]
\item Numerical and structural analysis of the representation gap
$\beta_n$ vs.\ $r_n$ across dimensions.
\item Collision-class enumeration to depth $30$, the structural
sequence battery, the $\sqrt{3}$-asymptotic of $\delta_d - 2 z_d$, and
the crossover $\delta_d > u_d$ for $d \ge 24$.
\item $\{3, 4, 5\}$-cascade observation for maximal outerplanar graphs
on $n \le 15$ vertices, the slack/denominator/precision analysis, and
the cascade-existence conjecture.
\item The standard 3-simplex remark and merged length-$10$-relation
cascades from Pythagorean triples sharing a common edge length.
\item Per-script hardware envelopes and runtime tables for the BFS
implementations.
\end{itemize}

These items are unchanged from earlier supplementary drafts and are
recorded here as a placeholder; consult the project repository for the
detailed records.

\section*{References}

References to literature are as in the main paper.  The additional
references invoked by the conditional routes above are:
\begin{itemize}[leftmargin=2em,itemsep=0.2em]
\item G.~Christol, \emph{Ensembles presque p\'eriodiques
$k$-reconnaissables}, TCS 9 (1979), 141--145.
\item J.-P.~Allouche, J.~Shallit, \emph{Automatic Sequences},
Cambridge University Press, 2003.
\item B.~Adamczewski, J.~P.~Bell, \emph{On the algebraic dependence of
$p$-automatic series and $G$-functions over function fields of
positive characteristic}, J.~Number Theory 180 (2017), 559--586.
\item J.~P.~Bell, M.~Coons, \emph{A Mahler dichotomy for analytic
functions of bounded height}, J.~Number Theory 137 (2014), 91--115.
\item B.~Adamczewski, J.~P.~Bell, \emph{A problem about Mahler
functions}, Annali Sc.~Norm.~Sup.~Pisa (5) 14 (2015), 541--562.
\item B.~Adamczewski, J.~P.~Bell, M.~Coons, \emph{A height gap theorem
for coefficients of Mahler functions}, J.~Eur.~Math.~Soc.\ 2023.
\item B.~Adamczewski, C.~Faverjon, \emph{Mahler's method in several
variables and finite automata}, Compositio Math.\ 156 (2020).
\item R.~P.~Stanley, \emph{Enumerative Combinatorics}, Vol.~2,
Cambridge University Press, 1999.  Thm.~6.4.6 (algebraic $\Rightarrow$
D-finite).
\item P.~Flajolet, R.~Sedgewick, \emph{Analytic Combinatorics},
Cambridge University Press, 2009.  Thm.~VII.7.
\item Y.~Antol\'in, L.~Ciobanu, \emph{Formal conjugacy growth series
and language complexity}, Math.\ Z.\ 285 (2017), no.~3, 1015--1041.
\end{itemize}

\end{document}
